\documentclass[10pt,a4paper]{amsart}
\usepackage[margin=19mm]{geometry}
\usepackage{amsmath,amssymb,mathtools}
\usepackage[scr=rsfs,cal=euler]{mathalfa}
\usepackage{xfrac}
\usepackage[unicode,hidelinks,bookmarks=false]{hyperref}
\newcommand{\ie}{i.e.\ }
\newcommand{\cf}{cf.\ }
\newcommand{\resp}{respectively\ }
\newcommand{\Subsection}{\subsection}
\providecommand{\nobreakdash}{\nobreak\hbox{-}}
\setlength{\emergencystretch}{2em}
\multlinegap0pt
\usepackage{bm}
\usepackage{bbm}
\usepackage{dsfont}
\usepackage{xspace}
\usepackage{cancel}
\usepackage{mathtools}
\usepackage[shortlabels]{enumitem}
\usepackage{amsthm}
\makeatletter
\@ifundefined{theorem}{\newtheorem{theorem}{Theorem}}{}
\@ifundefined{proposition}{\newtheorem{proposition}[theorem]{Proposition}}{}
\makeatother
\DeclareMathOperator{\GL}{GL}
\DeclareMathOperator{\Hom}{{Hom}}
\DeclareMathOperator{\ch}{char}
\DeclareMathOperator{\orb}{orb}
\DeclareMathOperator{\rank}{rank}
\newcommand{\VV}{\mathcal{V}}
\renewcommand{\k}{\Bbbk}
\title[Twisted homology jump loci, twisted Alexander polynomials, a]{Twisted homology jump loci, twisted Alexander polynomials, and $\Sigma$-invariants}
\author{Yongqiang Liu, Alexander I. Suciu}
\date{Source: arXiv:2605.28595v2 (CC BY 4.0)}
\begin{document}
\maketitle

\noindent\emph{Source and licence.} Twisted homology jump loci, twisted Alexander polynomials, and $\Sigma$-invariants. Version arXiv:2605.28595v2 (CC BY 4.0), distributed under the Creative Commons Attribution 4.0 International licence, as recorded in the arXiv licence metadata for this exact version and shown on the abstract page https://arxiv.org/abs/2605.28595v2. Full rights notice: licenses/arxiv-2605.28595-CC-BY-4.0.txt.\par\smallskip

\noindent\textit{Editorial scope.} Counted unit: the Proposition whose source label is \texttt{prop:compact-orbifold-group} in the article (source0.tex lines 2007--2030, source label \texttt{prop:compact-orbifold-group}), delivered together with the article's own complete original proof of it (source0.tex lines 2032--2111, one \texttt{proof} environment of 80 lines). No mathematics is written, repaired or re-derived: statement and proof are the article's own text line for line. Required context retained uncounted: none. Every definition, display and label of those lines is retained unchanged. Omitted from this package and not redistributed: 1--2006, 2031--2031, 2112--2937. Nothing inside a retained excerpt is omitted or altered: every retained body is delivered exactly as the article's lines are written, so no omission receipt of any kind accompanies this package. Citations: the retained excerpts carry no citation command at all, so no bibliography entry of the article is transcribed and no reference is redistributed. Reference adaptation: none was needed and none was made; every reference inside the delivered unit resolves inside it, so no unresolved reference or citation is produced. Printed slips kept literal: no printed word, symbol, hypothesis or number of the counted statement or of its proof was altered. Layout and class: the article's own class is \texttt{amsart} with packages including txfonts, amscd, amssymb, amsmath, amsbsy, amsthm, hyperref, tikz, enumitem, nccmath, relsize, xfrac, microtype, verbatim; the wrapper is the lane's pinned \texttt{amsart} editorial wrapper at 10pt on A4, which restates the theorem-like environments the retained text needs and adds the article's own macro definitions that the retained text needs (\texttt{\textbackslash GL}, \texttt{\textbackslash Hom}, \texttt{\textbackslash VV}, \texttt{\textbackslash ch}, \texttt{\textbackslash k}, \texttt{\textbackslash orb}, \texttt{\textbackslash rank}), with extra wrapper packages (bm, bbm, dsfont, xspace, cancel, mathtools, enumitem). The delivered numbering is the unit's own. Assets: no retained span contains a figure environment, a graphics-inclusion command, a PDF-page inclusion, a table or a TikZ picture; no image file is copied into this package, the package declares no asset dependency and no image file is redistributed with it. Licence: version arXiv:2605.28595v2 of this article is distributed under the Creative Commons Attribution 4.0 International licence, as recorded in the arXiv licence metadata for this exact version and shown on the abstract page https://arxiv.org/abs/2605.28595v2. A full rights notice is delivered at licenses/arxiv-2605.28595-CC-BY-4.0.txt. Characters: every retained excerpt is pure ASCII, so the delivered engine, XeLaTeX with its default Latin Modern text font, reports no missing character.
\begin{proposition}
\label{prop:compact-orbifold-group}
Let $\Gamma=\pi_1^{\orb}(C_g,\bm{\mu})$, and fix an algebraically
closed field $\k$. Let $\sigma \colon \Gamma \to \GL_r(\k)$ 
be a representation, sending generators $x_i,y_i,z_j$ to matrices
$X_i,Y_i,Z_j$, respectively, such that
\[
\prod_{i=1}^{g}[X_i,Y_i]\cdot \prod_{j=1}^{s} Z_j = I_r,
\qquad
Z_j^{\mu_j}=I_r \quad (1\le j\le s).
\]

Then the following hold:
\begin{enumerate}
\item If $g\ge 2$, or if the matrix $1+Z_j+\cdots+Z_j^{\mu_j-1}$ 
is not invertible for some $j$, then $\Delta^\sigma(\Gamma)=0$.

\item If $g=1$ and the matrices $1+Z_j+\cdots+Z_j^{\mu_j-1}$ 
are invertible for all $1\le j\le s$, then $\Delta^\sigma(\Gamma)=1$.
\end{enumerate}

In particular, the invertibility condition holds whenever
$Z_j=I_r$ and $\ch(\k)\nmid \mu_j$.
\end{proposition}    

\begin{proof}
Consider the compact orbifold group $\Gamma=\pi_1^{\orb}(C_{g},{\bm \mu})$ 
as above, and denote by $X=K(\Gamma,1)$ its classifying space. By Fox calculus, 
the boundary map $\partial_2$ in $C_*(\widetilde{X};\k)$ is given by the 
$(2g+s)\times (s+1)$ matrix 
\[ 
\begin{pmatrix} 
\label{matrix compact} 
\frac{z_{1}^{\mu_{1}}-1}{z_{1}-1} & 0 & \cdots & 0 & (\prod_{j=1}^s z_j)^{-1} \\
0 & \frac{z_{2}^{\mu_{2}}-1}{z_{2}-1} & \cdots & 0 & (\prod_{j=2}^s z_j)^{-1} \\
\vdots & \vdots & \vdots & \vdots & \vdots \\
0 & 0 & \cdots & \frac{z_{s}^{\mu_{s}}-1}{z_{s}-1} &  z_s^{-1} \\
0& 0 & \cdots & 0 & 1-x_1y_1x_1^{-1} \\
0& 0 & \cdots & 0 & x_1-[x_{1},y_{1}] \\
0& 0 & \cdots & 0 & [x_1,y_1] (1-x_2y_2x_2^{-1}) \\
0& 0 & \cdots & 0 & [x_1,y_1] (x_2-[x_{2},y_{2}]) \\
\vdots& \vdots & \vdots & \vdots & \vdots \\
0& 0 & \cdots & 0 & \prod_{k=1}^{g-1}[x_k,y_k] (1-x_g y_g x_g^{-1}) \\
0& 0 & \cdots & 0 & \prod_{k=1}^{g-1}[x_k,y_k] (x_g-[x_{g},y_{g}]) \\
\end{pmatrix}. 
\]
Here we adopt the notation $\frac{z_{j}^{\mu_{j}}-1}{z_{j}-1}\coloneqq 
1+z_j+\cdots+z_j^{\mu_j-1}$.

Consider a character $\rho\in \Hom^0(\Gamma, \k^{\times})$, 
the connected component of $\Hom(\Gamma,\k^{\times})$ 
containing the constant sheaf, which sends $x_i\to \rho_i$, 
$y_i\to \rho_{g+i}$, and $z_j\to 1$. 
Then $\partial_2^{\sigma \otimes \rho} $ has the following 
$(2g+s)\times (s+1)$ matrix:
\[ 
\begin{pmatrix} 
\label{matrix compact bis} 
\frac{ Z_{1}^{\mu_{1}}-I_r}{ Z_{1}-I_r} & 0 & \cdots & 0 & (\prod_{j=1}^s  Z_j)^{-1} \\
0 & \frac{ Z_{2}^{\mu_{2}}-I_r}{ Z_{2}-I_r} & \cdots & 0 & (\prod_{j=2}^s  Z_j)^{-1} \\
\vdots & \vdots & \vdots & \vdots & \vdots \\
0 & 0 & \cdots & \frac{Z_{s}^{\mu_{s}}-I_r}{ Z_{s}-I_r} & Z_s^{-1} \\
0& 0 & \cdots & 0 & I_r-\rho_{g+1}\cdot X_1Y_1X_1^{-1} \\
0& 0 & \cdots & 0 & \rho_1\cdot X_1-[X_{1},Y_{1}] \\
0& 0 & \cdots & 0 & [X_1,Y_1] (I_r-\rho_{g+2}\cdot X_2Y_2X_2^{-1}) \\
0& 0 & \cdots & 0 & [X_1,Y_1] (\rho_2\cdot X_2-[X_{2},Y_{2}]) \\
\vdots& \vdots & \vdots & \vdots & \vdots \\
0& 0 & \cdots & 0 & \prod_{k=1}^{g-1}[X_k,Y_k] (I_r-\rho_{2g}\cdot X_g Y_g X_g^{-1}) \\
0& 0 & \cdots & 0 & \prod_{k=1}^{g-1}[X_k,Y_k] (\rho_g\cdot X_g-[X_{g},Y_{g}])
\end{pmatrix},
\]
where $\frac{Z_{j}^{\mu_{j}}-I_r}{Z_{j}-I_r}\coloneqq 1+Z_j+\cdots+Z_j^{\mu_j-1}$.

Assume that for some $1\leq i\leq g$, one of the following holds:
\begin{itemize}[itemsep=2pt]
\item $\rho_i$ is not an eigenvalue of $[X_i,Y_i]X_i^{-1}$ and $X_i^{-1}$; or 
\item   $\rho_{g+i}^{-1}$ is not an eigenvalue 
of $X_iY_iX_i^{-1}$ and $Y_i$.
\end{itemize}
Then
\[
\rank \partial_2^{\sigma \otimes \rho} = r+ \sum_{j=1}^s 
\rank \tfrac{Z_{j}^{\mu_{j}}-I_r}{ Z_{j}-I_r}
\]
and $\rank \partial_1^{\sigma \otimes \rho} =r$. 
 Putting all together, we get that 
\begin{align*}
    \dim H_1(\Gamma; L_\sigma\otimes L_\rho)& 
    =r(2g+s)-\rank \partial_1^{\sigma \otimes \rho} 
    -\rank \partial_2^{\sigma \otimes \rho}\\
    &=r(2g-2+s)-\sum_{j=1}^s  \rank  \tfrac{Z_{j}^{\mu_{j}}-I_r}{ Z_{j}-I_r}.
\end{align*} 
 Hence if either $g\geq 2 $ or $ \frac{Z_{j}^{\mu_{j}}-I_r}{Z_{j}-I_r}$ is 
 not of full rank for some $j$, then 
 \[
 \VV_1(\Gamma,L_\sigma)\cap \Hom^0(\Gamma,\k^{\times})= 
 \Hom^0(\Gamma,\k^{\times}).
\]
In particular, the first Alexander polynomial $\Delta^\sigma(\Gamma)=0$ in this case.
On the other hand, if $g=1$ and $ \frac{Z_{j}^{\mu_{j}}-I_r}{ Z_{j}-I_r}$ 
has full rank  for all $1\leq j \leq s$, then 
$\VV_1(\Gamma,L_\sigma)\cap \Hom^0(\Gamma,\k^{\times})$ consists 
of only finitely many points. In particular, the first Alexander polynomial 
$\Delta^\sigma(\Gamma)=1$ in this case.
\end{proof}

\end{document}
